\chapter{Is Being Early an Advantage?}

Everyone has had the same small stab of regret:

\begin{quote}
Why did I only learn this now?
\end{quote}

Or, in its more painful form:

\begin{quote}
If only I had known earlier.
\end{quote}

The sentence feels natural, but it is not always useful. It often hides several different questions inside one emotion. Earlier than whom? Earlier in what condition? Earlier with what knowledge, what capital, what emotional strength, and what ability to act? If the same opportunity appeared again, would it really belong to you, or would you miss it again for a different reason?

Regret is not useless. It can be evidence. But it must be examined, not worshiped.

The common phrase ``first-mover advantage'' is dangerous for the same reason. It contains truth, but not enough truth. Sometimes early entry matters. Sometimes it matters enormously. In other situations, being early is simply another way to be wrong for a long time, to run out of capital before the market arrives, or to educate later competitors with one's own mistakes.

The point is not to deny the value of being early. The point is to put it in its proper place.

\section*{The Time Value of Knowledge}

In the previous chapter, we examined greed and temptation through the lens of return. Money has a time value. This is why interest exists. A unit of money today is not the same as a unit of money ten years from now, because time can be used to create more money.

Knowledge also has a time value.

This sounds obvious once stated, yet most people do not naturally think this way. Our languages have a common word for the time value of money: interest. We can argue about interest rates, compound interest, usury, debt, lending, and yield because the concept has a name.

But what is the common word for the time value of knowledge?

There is no ordinary word that carries the same function.

That absence matters. When a culture lacks a word, many people lack the mental handle. They may still experience the phenomenon, but they cannot easily inspect it. They feel regret after learning something late, but they do not see that the thing they missed was a form of knowledge interest.

Some knowledge produces returns only after it has been understood early enough and acted on for long enough. A person who learns probability before entering markets avoids errors that may destroy capital. A person who learns English early enough can access information, tools, and networks for decades. A person who learns how compounding works while still young has more years in which correct behavior can accumulate.

Knowledge can also produce negative interest.

A vague concept is not harmless. A wrong concept is not neutral. A person who believes that investing means short-term prediction will train the wrong reflex every day. A person who believes that all high yield is opportunity will misread risk every day. A person who believes that attention is cheap will squander the very asset needed for future freedom.

The longer a wrong concept remains in the mind, the more debt it creates.

\begin{center}
\begin{adjustbox}{max width=\linewidth}
\begin{tabular}{p{0.28\linewidth}p{0.27\linewidth}p{0.31\linewidth}}
\hline
Asset & Positive interest & Negative interest \\
\hline
Money & productive return over time & debt, fees, bad allocation \\
Knowledge & clearer action over time & repeated error from bad concepts \\
Attention & accumulated skill and judgment & scattered reaction and anxiety \\
\hline
\end{tabular}
\end{adjustbox}
\end{center}

This is why the careful definition of concepts has been repeated throughout this book. A clarified concept is not decoration. It is a repaired instrument. If the instrument is wrong, every measurement after it may be wrong.

\section*{Knowing Early Is Not Owning Early}

Some knowledge has direct commercial value. The earlier one knows it, the more valuable it may become, provided that one can act on it.

That last condition must not be omitted.

In early 2017, when Bitcoin again reached a historical high near \( \$1{,}200 \), many people suddenly said:

\begin{quote}
If only I had bought some Bitcoin earlier.
\end{quote}

A few days later, after a sharp fall, many of the same emotions appeared in another form. This pattern is ancient. When a price rises, people regret not buying. When it falls, they regret buying or feel relieved that they did not buy. Their emotional state follows the chart because their understanding has not moved ahead of the chart.

The key distinction is this:

\begin{quote}
Past correctness is not the same as future judgment.
\end{quote}

Looking backward, the value is visible. Everyone can see what already happened. If a person says, ``I should have bought then,'' he is often not expressing judgment. He is expressing hindsight. The real question is whether the value can continue from now into the future. If he cannot judge that, then his present self is not much different from his past self.

This is why some supposed opportunities never belonged to us, even if we once stood near them. If the same history were replayed, many people would still fail to buy, or would buy too little, or would sell too early, or would panic at the first major drawdown, or would be unable to withstand the social pressure of looking foolish for years.

To know early is one thing. To understand enough is another. To act is another. To hold through uncertainty is another still.

\section*{The Bitcoin Lesson}

Bitcoin began in 2009. Learning about it in 2011 was not actually first. By then, many people had known about it earlier. Some had mined it, discussed it, traded it, or ignored it. There was even the famous early physical transaction in which a large number of bitcoins bought a pizza. The person who received those coins did not thereby become a permanent winner. Knowing early did not guarantee keeping.

The deeper lesson is not ``be earlier next time.'' It is more demanding.

A real opportunity requires several conditions to meet at once:

\begin{enumerate}
\item You notice it early enough to matter.
\item You understand it deeply enough to avoid treating it as pure fantasy.
\item You have enough capital, but not so much that the opportunity looks too small to care about.
\item You keep learning after entering.
\item You grow fast enough as a person to carry the larger result later.
\end{enumerate}

This combination is rarer than beginners imagine.

There is a strange middle zone in capital. If one is too poor, one may not be able to allocate meaningful funds or withstand volatility. If one is already extremely rich, a small emerging asset may look irrelevant. In 2011, Bitcoin's total market value was small. Many brilliant people understood the idea quickly and still did not buy, not because they were foolish, but because the asset was too small relative to their world.

The right amount of wealth may be not too little and not too much. Enough to act, not so much that the scale blinds you.

Yet even this is not the largest point. The largest point is that early entry is not the core value when the target is long-term compound growth. Many people were early. Few held. Fewer held with understanding. Fewer still continued to improve fast enough to remain capable as the stakes grew.

A small position that becomes large tests a different person from the one who first bought it. If that person has not grown, the position may become unbearable.

\section*{Knowledge Interest Is Not Fixed}

Money interest can be calculated under fixed assumptions. If principal and rate are fixed, then more time produces more interest. In that setting, starting earlier gives an obvious advantage.

Knowledge is different. Its rate is not fixed. It depends heavily on the person.

A person can begin early and still accumulate confusion. He can read for years and become more rigid. He can invest for years and become more superstitious. He can work in an industry for a decade and learn only the habits of the average participant. In that case, time did not produce positive knowledge interest. It produced negative interest.

Another person may begin later but learn faster, ask better questions, correct errors sooner, and build a more reliable system. He is late in calendar time, but early in conceptual clarity.

The practical model is:

\[
\text{Result}
=
\text{Starting point}
\times
(1+\text{growth rate})^{\text{duration}}
\]

The starting point matters, but it is not everything. Growth rate and duration matter as well. For knowledge, growth rate is shaped by attention, humility, metacognition, feedback, reading, practice, and the willingness to revise oneself.

This is why one question should return again and again:

\begin{quote}
What is more important?
\end{quote}

Is it more important to have been early, or to be right for a long time? Is it more important to do something efficiently, or to do the correct thing? Is it more important to feel regret, or to find the exact reason one was unable to act?

In doing things, doing the right thing is more important than doing the thing right. If the direction is roughly correct and durable, even imperfect execution can accumulate. If the direction is wrong, efficient execution only makes the error larger.

\section*{Pressure Changes the Mind}

It is easy to imagine courage in the past. It is hard to exercise courage in the present.

Looking back, every line on the chart is settled. The price at which one ``should have bought'' and the price at which one ``should have sold'' are both visible. In memory, there is no uncertainty. In real time, every decision contains doubt, incomplete information, and pressure.

The size of the stake changes the mind.

A snooker player may produce perfect breaks in practice. In a major tournament, after reaching the final stages, the same player faces prize money, reputation, audience, and history. The physical task may be similar, but the psychological environment is different. Too much is at stake.

Investing works the same way. A person may say, ``I would have held.'' But would he have held when the amount became several years of income? Would he have held through a 50\% fall? Would he have held when intelligent friends mocked the asset? Would he have held when media declared it dead? Would he have held after becoming rich enough to feel that losing it would be unbearable?

Most self-confidence in investing is simply ignorance of future pressure.

This is not an insult. It is a warning. The market reveals what the imagination cannot test. A person should therefore be slow to say, ``I would have done better.'' The more useful statement is:

\begin{quote}
What ability was missing then, and how can I build it before the next opportunity?
\end{quote}

\section*{First-Mover Advantage and Later-Mover Advantage}

First-mover advantage is real in some situations. A company that enters early may gain users, data, brand, relationships, habits, distribution, patents, talent, or trust before competitors arrive. Networks can make this effect especially strong, because each additional participant can increase the value of the whole system.

A simple network table shows why growth can accelerate:

\begin{center}
\begin{adjustbox}{max width=\linewidth}
\begin{tabular}{lccccc}
\hline
Nodes & 2 & 3 & 4 & 5 & 6 \\
Connections & 1 & 3 & 6 & 10 & 15 \\
\hline
\end{tabular}
\end{adjustbox}
\end{center}

In a fully connected network, the number of possible connections is:

\[
\frac{n(n-1)}{2}
\]

This is one reason an early network can become hard to displace. The advantage is not merely that it appeared first. The advantage is that early appearance may lead to accumulated connections, and accumulated connections may create value that later entrants cannot easily copy.

But this is only one pattern. There is another.

The early entrant may arrive before the market is ready. It may spend heavily educating users. It may choose the wrong technology. It may build habits that later become liabilities. It may burn capital while later competitors wait, observe, and enter with a better product when demand is already proven.

Being early can be an advantage. It can also be the reason one does not survive until tomorrow.

A later entrant can have its own advantages:

\begin{enumerate}
\item The market has already been educated.
\item User demand is easier to observe.
\item Earlier mistakes are visible.
\item Technology and infrastructure may be cheaper.
\item The later entrant can focus on execution instead of proof of concept.
\end{enumerate}

Therefore the better question is not, ``Am I first?''

The better question is:

\begin{quote}
What kind of advantage exists here, and can I actually use it?
\end{quote}

\section*{The Venture Capital Illusion}

Risk capital invests in early-stage companies before they are publicly traded. Since the activity is early by design, many people assume that its returns must be generally superior to public-market investing.

The reality is more selective.

In the American venture capital industry, research has shown an extreme concentration of returns. A small minority of firms capture most of the industry's gains, and the leading group remains relatively stable over long periods. Remove the strongest funds, and the remaining industry record can look painfully ordinary or even poor.

This matters because ``early'' is often confused with ``profitable.'' Early-stage access does not automatically create superior return. It may simply expose the investor to earlier failure.

Consider Facebook. Before it became public, the best early opportunities belonged to a narrow group of investors with unusual access, judgment, networks, and capital. That was not an ordinary person's opportunity. Wanting it as if it naturally belonged to us is a form of greed.

After Facebook became public, ordinary investors could buy its shares. Was that too late? Not necessarily. A person who bought after the initial public offering and held through the following years could still have done very well. The opportunity was no longer private, no longer earliest, and no longer reserved for insiders. Yet it remained meaningful.

This example is valuable because it breaks a common illusion:

\begin{quote}
If I was not early, I have already missed it.
\end{quote}

Not always. Sometimes the public, later, more visible opportunity is the one that actually belongs to you, because it is the one you can access, study, buy, and hold within your own competence.

The investor's task is not to envy Peter Thiel's opportunity. It is to recognize one's own.

\section*{When a Tool Arrives Late Enough}

At the end of 2016, a small team was asked to build a paid lecture tool inside WeChat. The team resisted. Their argument was simple: too many similar tools already existed, so there was no first-mover advantage.

The tool was built anyway. Within a month, it had accumulated hundreds of thousands of paying users, tens of thousands of daily active users, and millions of yuan in transaction volume.

The lesson is not that first-mover advantage never matters. The lesson is that the phrase can become a lazy substitute for judgment. In that case, the team saw existing competitors and stopped thinking. They did not ask whether the demand was strong enough, whether distribution was available, whether execution could be faster, whether trust already existed, whether timing had improved, or whether earlier products had arrived too early.

A concept becomes dangerous when it turns into a hammer. After learning the phrase ``first-mover advantage,'' some people begin to see every situation as a contest to be first. But advantage has many forms. Timing is one. Distribution is one. Trust is one. Capital is one. Product quality is one. Persistence is one. Existing audience is one. Lower cost from later technology is one.

A mature person does not ask only whether he is early. He asks whether he has a usable advantage.

\section*{A Better Review of Missed Opportunities}

Regret should be converted into review. The review should be concrete.

When you feel, ``I learned too late,'' ask:

\begin{enumerate}
\item What exactly did I learn late?
\item Was the opportunity visible to people like me at the time?
\item If I had seen it, would I have understood it?
\item If I had understood it, would I have acted?
\item If I had acted, would I have held through pressure?
\item What missing concept, skill, capital structure, or habit prevented me?
\end{enumerate}

This review often reduces unnecessary self-blame. Some opportunities truly did not belong to you then. They required access you did not have, knowledge you had not built, or psychological strength you had not trained. Recognizing this is not an excuse. It is the beginning of accurate preparation.

The same review may also reveal avoidable errors. Perhaps you did not read enough. Perhaps you dismissed something because it sounded strange. Perhaps your English was too weak to follow the primary information. Perhaps you had no savings because spending habits consumed all optionality. Perhaps you were too addicted to security to tolerate uncertainty. Perhaps your attention was scattered when it should have been studying.

The goal is not to suffer more. The goal is to become more precise.

A missed opportunity can teach only if the reason for missing it is named.

\section*{The Luck of Some Shortcomings}

Not every missed opportunity is a tragedy. Sometimes a weakness protects you.

A person with a poor car may drive more slowly and avoid an accident. A person without enough capital may be unable to enter a fraudulent scheme. A person too cautious to chase a bubble may later discover that caution saved him. A person who lacks access to a fashionable private deal may avoid a loss that insiders themselves did not understand.

This should not become pride in weakness. Weakness is still weakness. But it should soften the habit of blaming fate. Some disadvantages are expensive. Some are strangely lucky. Life is often more mixed than the regretful mind admits.

The better attitude is calm inventory. What has hurt me? What has protected me? Which protections should be replaced by real ability? Which weaknesses must be corrected before they stop being lucky and become permanent limits?

A person who can think this way becomes less dramatic and more useful to himself.

\section*{Do Not Use This as an Excuse}

There is one final danger.

After hearing that being early is not decisive, a person may say:

\begin{quote}
Then I do not need to start yet.
\end{quote}

This is the old demand for excuses in a new form. Earlier chapters have shown how powerful excuse-making can be. The mind that wants delay can use any idea as raw material. If early matters, it says, ``I am already late, so why act?'' If early does not always matter, it says, ``Then there is no urgency, so why act?''

Both are evasions.

The correct conclusion is:

\begin{quote}
Being early is not the only advantage, but never starting is always a disadvantage.
\end{quote}

Action creates feedback. Feedback refines concepts. Refined concepts improve judgment. Improved judgment raises the interest rate of knowledge. Without action, this loop does not begin.

\begin{center}
\begin{adjustbox}{max width=\linewidth}
\begin{tikzpicture}[font=\scriptsize, >=stealth, node distance=.95cm]
\node[draw, rounded corners=2pt, align=center, minimum width=2.35cm] (act) {Act};
\node[draw, rounded corners=2pt, align=center, right=of act, minimum width=2.35cm] (feed) {Feedback};
\node[draw, rounded corners=2pt, align=center, right=of feed, minimum width=2.35cm] (concept) {Clearer\\concepts};
\node[draw, rounded corners=2pt, align=center, below=of feed, minimum width=2.35cm] (judge) {Better\\judgment};
\draw[->] (act) -- (feed);
\draw[->] (feed) -- (concept);
\draw[->] (concept) |- (judge);
\draw[->] (judge) -| (act);
\end{tikzpicture}
\end{adjustbox}
\end{center}

Learning English, exercising, training metacognition, saving capital, building a small investment record, writing regularly, studying probability, and improving execution are all correct things whose value does not depend on being first. They are not late because someone else began earlier. Their returns depend on beginning, continuing, and improving.

The practical conclusion is plain. Stop worshiping the fantasy of an earlier self. Accept the cold fact that you are learning this now. Then ask what can still compound from here.

First is not enough. Late is not fatal. Wrong is expensive. Not starting is worse.
